\begin{table}[!htbp]
\par\noindent\begin{minipage}{\linewidth}
\small
\setlength{\tabcolsep}{4pt}
\renewcommand{\arraystretch}{1.13}
\captionof{table}{Modelos de embeddings y reglas principales de representación.}\label{tab:T02}
\begin{tabular}{@{}>{\raggedright\arraybackslash}m{3.3cm}>{\centering\arraybackslash}m{1.6cm}>{\centering\arraybackslash}m{1.65cm}>{\raggedright\arraybackslash}m{3.0cm}>{\raggedright\arraybackslash}m{4.9cm}@{}}
\toprule
\textbf{Modelo} & \textbf{Tamaño} & \textbf{Límite de tokens} & \textbf{Vector del artículo} & \textbf{Referencia} \\
\midrule
\textbf{SPECTER} & 768 & 512 & CLS & \citep{cohan-etal-2020-specter} \\
\addlinespace[2pt]
\textbf{SPECTER2} & 768 & 512 & CLS + adaptador & \citep{singh-etal-2023-scirepeval} \\
\addlinespace[2pt]
\textbf{SciNCL} & 768 & 512 & CLS & \citep{ostendorff-etal-2022-neighborhood} \\
\addlinespace[2pt]
\textbf{SciBERT} & 768 & 512 & Media & \citep{beltagy-etal-2019-scibert} \\
\addlinespace[2pt]
\textbf{BERT} & 768 & 512 & Media & \citep{devlin-etal-2019-bert} \\
\addlinespace[2pt]
\textbf{MPNet} & 768 & 384 & Media & \citep{reimers-gurevych-2019-sentence} \\
\addlinespace[2pt]
\textbf{MiniLM} & 384 & 256 & Media & \citep{reimers-gurevych-2019-sentence} \\
\addlinespace[2pt]
\textbf{PubMedBERT / BiomedBERT} & 768 & 512 & Media & \citep{reff9ebfb9307} \\
\addlinespace[2pt]
\textbf{BioBERT} & 768 & 512 & Media & \citep{ref4f3e28932e} \\
\addlinespace[2pt]
\textbf{SimCSE} & 768 & 512 & CLS & \citep{gao-etal-2021-simcse} \\
\bottomrule
\end{tabular}
\par\smallskip{\small Tamaño indica cuántas coordenadas contiene el vector, no cuántas palabras lee el modelo. MPNet y MiniLM utilizan all-mpnet-base-v2 y all-MiniLM-L6-v2, respectivamente. Los límites cuentan tokens, las piezas en que el modelo divide el texto, incluidos sus símbolos especiales. Media promedia las salidas finales de tokens (la última capa), con símbolos especiales y sin posiciones vacías de relleno. CLS usa la salida de la primera posición antes de la transformación adicional llamada pooler. SPECTER2 utiliza también su adaptador de proximidad, un componente adicional entrenado. En MPNet/MiniLM, la referencia describe el marco Sentence-BERT, no el entrenamiento exacto de esas versiones. Las reglas alternativas no son modelos adicionales. Versiones, formatos y recorte: Tabla S2.\par}
\end{minipage}\par
\end{table}
